\documentclass[aspectratio=169]{beamer}
\usepackage{spc}
\usepackage{textpos}
\begin{document}

\begin{frame}
  \title{\vspace{-5ex}\darkblue Previous YFT Assessment from 2023\\[2ex]
    \it\large\darkgray Results, Innovations, Issues}
  \author{\vspace{-10ex}\darkgray\bf Arni Magnusson}
  \date{\darkgreen SPC PAW Meeting, Nouméa\\[0.5ex]
    24 March 2026}
  \titlepage
  \begin{textblock}{0}(10,0.5)
    \includegraphics[width=4cm]{fish_yellowfin}
  \end{textblock}
\end{frame}

% ______________________________________________________________________________

\begin{frame}{Overview}
  \vspace{2ex}
  \begin{itemize}
    \item[] {\bf\darkblue Data and Model Fit}
    \comment{}\\[4ex]
    \item[] {\bf\darkblue Main Results}
    \comment{reference points, stock status}\\[4ex]
    \item[] {\bf\darkblue Innovations}
    \comment{simplified regional structure, Lorenzen M,\\[1ex]
      \h{15.4ex} catch-conditioned, estimation uncertainty}\\[4ex]
    \item[] {\bf\darkblue Issues}
    \comment{convergence/jittering, narrow confidence intervals,\\[1ex]
      \h{9.1ex} likelihood conflict between length and weight comps}\\[4ex]
  \end{itemize}
\end{frame}

% ______________________________________________________________________________

\begin{frame}\Large
  \centering\darkgreen\bf Innovations in 2023 Assessment
\end{frame}

% ______________________________________________________________________________

\begin{frame}{Simplified regional structure}
  \includegraphics[width=0.47\textwidth]{previous_9_regions}
  \raisebox{13ex}{$\Rightarrow$}
  \includegraphics[width=0.47\textwidth]{previous_5_regions}
\end{frame}

% ______________________________________________________________________________

\begin{frame}{Lorenzen natural mortality}
  \vspace{4ex}
  \centering\includegraphics[height=0.6\textheight]{natmort}
\end{frame}

% ______________________________________________________________________________

\begin{frame}{Catch-conditioned method}\small
  MULTIFAN-CL User Manual, section 4.5:\\[3ex]
  \includegraphics[width=\textwidth]{catch_conditioned}
\end{frame}

% ______________________________________________________________________________

\begin{frame}{Estimation uncertainty}
  \vspace{1ex}
  \centering\includegraphics[height=0.7\textheight]{estimation_uncertainty}
\end{frame}

% ______________________________________________________________________________

\begin{frame}\Large
  \centering\darkgreen\bf Issues in 2023 Assessment
\end{frame}

% ______________________________________________________________________________

\begin{frame}{Convergence and jittering}
  \vspace{1ex}
  \centering\includegraphics[height=0.65\textheight]{jittering}
\end{frame}

% ______________________________________________________________________________

\begin{frame}{Narrow confidence intervals}
  \begin{itemize}
    \item Incorporating estimation uncertainty had no real effect on the overall
    uncertainty.\\[0.5ex]
    For example, the confidence interval for
    $SB_\mathrm{recent}/SB_{F=0}$:\\[1ex]
    \quad $0.42\!-\!0.52$ without estimation uncertainty\\[0.5ex]
    \quad $0.42\!-\!0.52$ including estimation uncertainty\\[6ex]
    \item Also, the confidence interval around the estimated M is quite narrow:
    $0.11\!-\!0.13$\\[1ex]
  \end{itemize}
  \centering\includegraphics[width=0.89\textwidth]{natmort_confint_2023}
  \vspace{2ex}
\end{frame}

% ______________________________________________________________________________

\begin{frame}{Likelihood conflict between length and weight comps}
  \vspace{1ex}
  \centering\includegraphics[height=0.75\textheight]{likelihood_profile}
\end{frame}

% ______________________________________________________________________________

\begin{frame}{Summary}
  \vspace{2ex}
  \begin{itemize}
    \item[] {\bf\darkblue Innovations in 2023 Assessment}
    \comment{simplified regional structure, Lorenzen M,\\[1ex]
      \h{38.5ex} catch-conditioned, estimation uncertainty}\\[8ex]
    \item[] {\bf\darkblue Issues in 2023 Assessment}
    \comment{convergence/jittering, narrow confidence intervals,\\[1ex]
      \h{32.2ex} likelihood conflict between length and weight comps}\\[4ex]
  \end{itemize}
\end{frame}

% ______________________________________________________________________________

\begin{frame}
  \vspace{4ex}
  These slides:\\[2ex]
  \begin{center}
    \blue\url{https://github.com/PacificCommunity/ofp-sam-yft-2026-assessment}
  \end{center}
\end{frame}

\end{document}
